\pagenumbering{arabic}
\pagestyle{fancy}

% فصل ۱: کلیات
\chapter{کلیات}

\section{مقدمه}

در حوزه مهندسی سیستم‌های نرم‌افزاری معاصر، پیچیدگی به‌عنوان یکی از بارزترین ویژگی‌ها خودنمایی می‌کند. محیط عملیاتی پویا نقشی محوری در شکل‌دهی به این پیچیدگی دارد و سیستم‌ها را ناگزیر می‌سازد تا با شرایط نامشخصی که اغلب تنها پس از عملیاتی شدن بروز می‌یابند، مقابله کنند.

سیستم‌های سنتی ایستا یا مبتنی بر قواعد، انعطاف‌پذیری لازم برای پاسخ به چنین تغییراتی در زمان اجرا را ندارند. در مقابل، سیستم‌های خودوفقی (\lr{SASs}) مجهز به حلقه‌های بازخوردی هستند که به آن‌ها اجازه می‌دهد تا محیط و وضعیت خود سیستم را پایش کنند، انحرافات از رفتار مورد انتظار را تحلیل نمایند و به‌طور خودمختار استراتژی‌ها و تغییرات وفقی را برنامه‌ریزی و اجرا کنند.

تمرکز این رساله بر سیستم‌هایی با فضای وفقی گسسته و بزرگ است که نیاز فوری به راهکارهایی کارآمد برای انتخاب گزینه‌های بهینه در زمان اجرا دارند.

\section{اهداف رساله}

این رساله در راستای تکمیل شکاف‌های تحقیقاتی پیشین \cite{RS_DRL_Article} و پاسخ به چالش‌های حل‌نشده، یک هدف اصلی و سه هدف فرعی را دنبال می‌کند:

\subsection*{هدف اصلی}

\textbf{توسعه و اعتبارسنجی روش یادگیری تقویتی عمیق با شکل‌دهی تصادفی پاداش (\lr{RS-DRL}) به‌گونه‌ای که بر مشکل بنیادین گیرافتادن در بهینه‌های محلی در فضاهای وفقی گسسته و بزرگ غلبه کند. این روش سیستم‌های خودوفقی را قادر می‌سازد تا بدون نیاز به بازآموزی پرهزینه و اخلال‌گر، عملکردی نزدیک به بهینه را به‌طور پیوسته در شرایط عدم قطعیت محیطی حفظ نمایند.}

این هدف اصلی، هسته‌ی نوآوری رساله (مکانیزم شکل‌دهی تصادفی پاداش) را مستقیماً به چالش اساسی شناسایی‌شده در بیان مسئله پیوند می‌دهد.

\subsection*{اهداف فرعی}

در راستای تحقق هدف اصلی، اهداف فرعی زیر تعریف شده‌اند:

\begin{enumerate}
    \item \textbf{طراحی معماری مقیاس‌پذیر و یکپارچه:} ایجاد معماری مدولاری که عامل \lr{RS-DRL} را به‌طور ناهمگام با حلقه‌ی خودمختار \lr{MAPE-K} ادغام نموده و امکان اعتبارسنجی صوری تصمیمات پیش از اجرا را فراهم آورد. این هدف تضمین می‌کند که راهکار پیشنهادی نه تنها هوشمند، بلکه قابل‌اعتماد و قابل‌استقرار در سیستم‌های زمان‌اجرای واقعی باشد.
    
    \item \textbf{مدیریت کارآمد فضاهای وفقی گسسته با ابعاد بسیار بزرگ:} نشان دادن این که روش \lr{RS-DRL} قادر است در سیستم‌هایی با فضای وفقی تا ۴۰۹۶ گزینه (مانند \lr{DeltaIoTv2})، بدون افت محسوس عملکرد، گزینه‌های بهینه یا نزدیک به بهینه را در زمان کوتاه انتخاب کند. این هدف، مقیاس‌پذیری عملی راهکار را برای کاربردهای دنیای واقعی اثبات می‌نماید.
    
    \item \textbf{ارزیابی تجربی جامع و اثبات برتری آماری:} مقایسه‌ی کمی دقیق روش \lr{RS-DRL} با روش‌های پیشرفته‌ی موجود (شامل \lr{$\epsilon$-greedy DQN}، \lr{HER}، \lr{Soft HER}، \lr{DLASER} و \lr{ML4EAS}) از نظر عملکرد مجانبی، سرعت همگرایی، و بهینه‌سازی چندهدفه، همراه با ارائه‌ی تحلیل‌های آماری معنی‌دار (مانند \lr{t-test} و فواصل اطمینان) برای تصدیق برتری روش پیشنهادی.
\end{enumerate}


\section{چالش‌ها}

چالش‌ها و منابع عدم قطعیت در سیستم‌های مختلف به شکل‌های گوناگونی بروز می‌یابند. به‌عنوان مثال، در سیستم‌های مستقر در محیط‌های پویا نظیر برنامه‌های اینترنت اشیا (\lr{IoT}) در شهرهای هوشمند، تغییرات آب و هوایی تأثیر قابل توجهی می‌گذارد. همچنین محیط‌های ابری با تقاضای منابع متغیر، امنیت سایبری با تهدیدات نوظهور و مداوم، تجارت الکترونیک و شبکه‌های اجتماعی با رفتار غیرقابل پیش‌بینی کاربران، و سیستم‌های توزیع‌شده با تأخیر شبکه، خرابی سخت‌افزار و تغییرات زیرساختی روبه‌رو هستند.

هنگامی که سیستم با فضای وفقی گسترده و گزینه‌های متعدد تغییر مواجه می‌شود، انتخاب گزینه‌ای وفقی که با اهداف سیستم همسو باشد، از نظر محاسباتی بسیار پرهزینه و دشوار، و در برخی موارد تقریباً غیرممکن می‌گردد.

روش‌های موجود برای مدیریت عدم قطعیت در سیستم‌های خودوفقی با محدودیت‌های قابل‌توجهی مواجه هستند. به‌عنوان مثال، روش‌های مبتنی بر یادگیری ماشین نیازمند داده‌های آموزشی برچسب‌دار هستند که نمایانگر محیط سیستم باشند. به دست آوردن چنین داده‌هایی در زمان طراحی به دلیل عدم قطعیت فراوان بسیار دشوار است. از طرفی، روش‌های مبتنی بر جستجو (مانند الگوریتم‌های صعود تپه و الگوریتم‌های ژنتیک) با مشکلات مقیاس‌پذیری مواجه هستند و عملکرد آن‌ها با تغییر شرایط محیطی دچار افت محسوس می‌شود.

یادگیری تقویتی عمیق (\lr{DRL}) علیرغم قدرت بالا، خود با چالش‌های اساسی نظیر گیرافتادن در بهینه‌های محلی مواجه است. این چالش در سیستم‌هایی با فضای وفقی گسسته و بزرگ به‌مراتب تشدید می‌شود؛ زیرا هرچه تعداد گزینه‌های وفقی بیشتر باشد، احتمال همگرایی مدل به راه‌حل‌های بهینه‌ی محلی به جای بهینه‌های سراسری افزایش می‌یابد و کاوش کافی در فضای عمل دشوارتر می‌گردد. بهینه‌های محلی به راه‌حل‌هایی اشاره دارند که تنها در یک منطقه محدود از فضای مسئله بهینه هستند و به صورت سراسری بهترین راه‌حل ممکن محسوب نمی‌شوند. در زمینه \lr{DRL}، این بدان معناست که مدل ممکن است سیاستی را یاد بگیرد که تحت شرایط خاص به خوبی عملکرد مناسبی دارد، اما در صورت تغییر محیط قادر به تعمیم‌پذیری و حفظ عملکرد مطلوب نیست.

علاوه بر این، آموزش مجدد یا بازآموزی (\lr{Retraining}) مدل در زمان اجرا فرآیندی پرهزینه است. بازآموزی شامل مراحل متعددی مانند جمع‌آوری داده‌های جدید، ارزیابی مجدد ساختار و پارامترهای مدل، و تست مدل به‌روزشده برای اطمینان از عملکرد بهتر آن نسبت به نسخه قبلی است. در طول این فرآیند، سیستم ممکن است به صورت غیربهینه عمل کند که در سیستم‌های حیاتی، این دوره عملکرد نامطلوب می‌تواند پیامدهای جبران‌ناپذیری به همراه داشته باشد.

\section{بیان مسئله}

سیستم‌های خودوفقی مدرن در محیط‌های پویا و غیرقابل‌پیش‌بینی نظیر شبکه‌های اینترنت اشیاء، رایانش ابری و شهرهای هوشمند عمل می‌کنند. این سیستم‌ها اغلب با فضاهای وفقی گسسته و بزرگ مواجه هستند که می‌تواند شامل صد‌ها یا هزاران گزینه پیکربندی مختلف باشد (برای مثال، شبکه حسگر \lr{DeltaIoTv2} دارای ۴۰۹۶ گزینه وفقی است). 

مدل‌های سنتی یادگیری تقویتی عمیق (\lr{DRL}) که برای تصمیم‌گیری در این فضاها به کار گرفته می‌شوند، با دو چالش بنیادین مواجه‌اند:

\begin{enumerate}
    \item \textbf{گیرافتادن در بهینه‌های محلی:} الگوریتم‌هایی مانند \lr{$\epsilon$-greedy DQN} به دلیل ماهیت کاوشگری محدود خود، اغلب در نقاط بهینه محلی متوقف می‌شوند و قادر به کشف سیاست‌های سراسری بهینه نیستند \cite{du2023alleviating}. این مشکل در فضاهای وفقی بزرگ به مراتب حادتر می‌شود.
    
    \item \textbf{نیاز به بازآموزی مکرر و پرهزینه:} در مواجهه با تغییرات محیطی (رانش مفهومی \cite{gheibi2024dealing})، مدل‌های سنتی نیاز به بازآموزی کامل دارند. این فرآیند زمان‌بر، محاسباتی سنگین است و در طول دوره بازآموزی، سیستم عملکرد نامطلوبی از خود نشان می‌دهد که در سیستم‌های حیاتی می‌تواند پیامدهای جبران‌ناپذیری به همراه داشته باشد.
\end{enumerate}

روش‌های پیشرفته‌تری مانند \lr{HER} و \lr{Soft HER} نیز اگرچه در حوزه محیط‌های هدف‌محور موفق بوده‌اند، اما برای سیستم‌های خودوفقی با اهداف مبهم و آستانه‌ای (مانند بهینه‌سازی همزمان مصرف انرژی، تأخیر و تلفات بسته) طراحی نشده‌اند و نمی‌توانند بدون دانش دامنه خاص عمل کنند.

بنابراین، مسئله اصلی این رساله را می‌توان به صورت زیر فرموله کرد:

\begin{center}
\textbf{طراحی روش یادگیری تقویتی عمیق که (۱) بدون نیاز به بازآموزی مکرر و پرهزینه، از بهینه‌های محلی در فضاهای وفقی گسسته و بزرگ (تا ۴۰۹۶ گزینه) عبور کند، (۲) قابلیت تعمیم‌پذیری به شرایط محیطی متغیر را داشته باشد، و (۳) مستقل از دانش دامنه خاص عمل کند.}
\end{center}

\section{سوالات تحقیق}

با توجه به محدودیت‌های رویکردهای پیشین و نیز دستاوردهای آن‌ها که در بخش چالش‌ها تبیین گردید \cite{RS_DRL_Article}، این رساله در صدد پاسخ به پرسش‌های اساسی زیر است:

\begin{enumerate}
  \item \textbf{سوال اول:} چگونه می‌توان روشی مبتنی بر یادگیری تقویتی عمیق طراحی کرد که به‌طور مؤثر از بهینه‌های محلی در فضاهای وفقی گسسته و بزرگ، بدون نیاز به بازآموزی مکرر و پرهزینه، عبور کند؟
  \item \textbf{سوال دوم:} چگونه می‌توان مکانیزم‌های شکل‌دهی تصادفی پاداش (\lr{Random Reward Shaping}) را با بازپخش تجربه (\lr{Experience Replay}) ادغام کرد تا تعمیم‌پذیری سیستم در تغییرات شرایط محیطی بهبود یابد؟
  \item \textbf{سوال سوم:} چگونه روش پیشنهادی می‌تواند فضاهای وفقی گسسته و بزرگ (از ۲۱۶ تا ۴۰۹۶ گزینه‌) را در حین حفظ عملکردی نزدیک به بهینه، به شکل کارآمدی مدیریت کند؟
  \item \textbf{سوال چهارم:} تا چه حد روش پیشنهادی (\lr{RS-DRL}) از روش‌های موجود نظیر یادگیری ای‌گرایانه حریصانه (\lr{$\epsilon$-greedy DQN})، \lr{HER}، \lr{Soft HER}، \lr{DLASER} و \lr{ML4EAS} از منظر عملکرد مجانبی، سرعت همگرایی، و بهینه‌سازی چندهدفه، برتری و عملکرد بهتری از خود نشان می‌دهد؟
\end{enumerate}

\section{روش پیشنهادی}

\subsection{مرور کلی بر راهکار \lr{RS-DRL}}
راهکار پیشنهادی در این رساله مبتنی بر توسعه یک معماری یادگیری هوشمند است که به جای تکیه بر دانش ایستا، از طریق تعامل مستقیم با محیط، سیاست‌های بهینه را استخراج می‌کند. این روش مسئله تطبیق را به‌عنوان یک فرآیند تصمیم‌گیری مارکوف (\lr{MDP}) مدل‌سازی کرده و با ادغام مدل‌های عمیق و لایه مدیریت خودمختار، یک چرخه پیوسته از یادگیری و اجرا را فراهم می‌آورد. هسته نوآورانه این روش، مکانیزم «شکل‌دهی تصادفی پاداش» است که با بازطراحی تصادفی برخی تجربیات ناموفق گذشته، عامل را قادر به فرار از بهینه‌های محلی بدون نیاز به بازآموزی پرهزینه می‌سازد. تمرکز اصلی بر ارائه چارچوبی است که در آن، عامل یادگیرنده نه تنها بر اساس پاداش‌های آنی، بلکه با هدف حفظ عملکرد مطلوب در بازه زمانی طولانی، کاهش نیاز به بازآموزی‌های مکرر و پایداری معیارهای کیفی سیستم نظیر مصرف انرژی و زمان تأخیر، عمل کند.

\subsection{یکپارچه‌سازی با حلقه کنترلی \lr{MAPE-K}}
روش پیشنهادی با بهره‌گیری از چارچوب استاندارد \lr{MAPE-K}، فرآیند تطبیق را به صورت ساختاریافته مدیریت می‌کند. در این معماری، عامل \lr{RS-DRL} در مؤلفه برنامه‌ریزی (\lr{Plan}) مستقر شده و خروجی‌های آن (گزینه وفقی پیشنهادی) پیش از اجرا توسط ابزار تأیید صوری \lr{UPPAAL-SMC} از نظر ایمنی و قابلیت اطمینان ارزیابی می‌گردند. سپس گزینه‌های تأییدشده بر روی سیستم مدیریت‌شونده اعمال می‌شوند. این یکپارچگی تضمین می‌کند که تصمیمات هوشمندانه‌ی عامل \lr{DRL} تنها پس از اعتبارسنجی صوری به اجرا درآیند.

\section{نوآوری‌های اصلی رساله}

بنیانِ علمی‌و سهم پژوهشی این پایان‌نامه در قالب سه نوآوری اصلی زیر تبیین می‌گردد که هدف نهایی آن‌ها بهبود همگرایی و تعمیم‌پذیری در سیستم‌های خودوفقی است:

\begin{enumerate}
  \item \textbf{نوآوری اول: الگوریتم یادگیری تقویتی با شکل‌دهی تصادفی پاداش (\lr{RS-DRL}):} 
  ارائه یک مکانیزم نوین که برخلاف روش‌های سنتی، برخی از تجربیات ناموفق گذشته را به صورت تصادفی بازطراحی می‌کند. با تخصیص پاداش مثبت به این تجربیات، فرآیند یادگیری تسریع یافته و از همگرایی مدل به بهینه‌های محلی جلوگیری می‌شود؛ بدین ترتیب، مدل بدون نیاز به بازآموزی مکرر و پرهزینه در برابر تغییرات محیطی مقاوم می‌گردد. این رویکرد الهام‌گرفته از یادگیری انسانی از اشتباهات است و بدون نیاز به دانش دامنه خاص، تعمیم‌پذیری مدل را در محیط‌های متنوع افزایش می‌دهد.
  
  \item \textbf{نوآوری دوم: ادغام معماری \lr{RS-DRL} در حلقه \lr{MAPE-K}:} 
  طراحی یک معماری تلفیقی که در آن فرآیند آموزش و استنتاج مدل عمیق به صورت ناهمگام با حلقه کنترل استاندارد سیستم‌های خودوفقی هماهنگ شده است. برخلاف رویکردهای پیشین که یادگیری را تنها به‌صورت همزمان یا آفلاین انجام می‌دهند، این معماری اجازه می‌دهد تا دانش آموخته شده در حافظه بازپخش به‌طور پیوسته در اتخاذ تصمیمات وفقی زمان اجرا نقش ایفا کند.
  
  \item \textbf{نوآوری سوم: مدیریت فضای وفقی مقیاس‌پذیر و گسسته:} 
  ارائهی راهکاری کارآمد برای هدایت و انتخاب پیکربندی در سیستم‌هایی با ابعاد بسیار بزرگ (تا ۴۰۹۶ گزینه‌ی وفقی). در حالی که روش‌های سنتی مبتنی بر جستجو (مانند الگوریتم‌های ژنتیک و صعود تپه) در این ابعاد با چالش نفرین ابعاد و افت کارایی مواجه می‌گردند، این روش نشان داده است که می‌تواند با حفظ پایداری و بهبود چشم‌گیر معیارهایی نظیر زمان تأخیر و هدررفت داده‌ها، بر این چالش‌ها در محیط‌های پویای اینترنت اشیاء غلبه نماید.
\end{enumerate}

\section{مروری بر فصول رساله}

ساختار رساله به گونه‌ای طراحی شده است که به‌طور نظام‌مند روند توسعه، اجرا و ارزیابی رویکرد پیشنهادی را در اختیار خواننده قرار دهد. در ادامه، مرور خلاصه‌ای از هر یک از فصول آورده شده است:

\textbf{فصل ۱: کلیات}

در این فصل، کلیات پژوهش، بیان مسئله، اهداف و نوآوری‌های رساله به همراه سوالات تحقیق و مروری بر ساختار کل رساله ارائه می‌گردد.

\textbf{فصل ۲: مبانی نظری و مدل سیستم}

در این فصل، مفاهیم بنیادین نظری و مدل صوری سیستم ارائه می‌گردد. این فصل شامل مبانی سیستم‌های خودوفقی (حلقه \lr{MAPE-K})، منابع عدم قطعیت و رویکردهای مدل‌سازی زمان اجرا (از جمله ابزار \lr{UPPAAL-SMC})، و اصول یادگیری تقویتی (فرآیند تصمیم‌گیری مارکوف و الگوریتم \lr{DQN}) است. سپس، سیستم شبکه‌ی \lr{DeltaIoT} به‌عنوان مطالعه‌ی موردی معرفی شده و فضای حالت، فضای عمل و توابع پاداش به صورت ریاضی فرمول‌بندی می‌شوند. در نهایت، مسئله بهینه‌سازی سیاست وفقی به صورت صوری بیان می‌گردد.

\textbf{فصل ۳: مرور ادبیات و کارهای مرتبط}

با بهره‌گیری از چارچوب نظری فصل دوم، این فصل یک بررسی انتقادی و جامع از رویکردهای موجود برای مدیریت عدم قطعیت در سیستم‌های خودوفقی ارائه می‌دهد. روش‌های پیشین در سه دسته اصلی بررسی شده و با شناسایی شکاف تحقیقاتی، ضرورت ارائه‌ی راهکار نوین \lr{RS-DRL} توجیه می‌گردد.

\textbf{فصل ۴: روش پیشنهادی \lr{RS-DRL}: معماری و الگوریتم}

این فصل به تشریح دقیق الگوریتم \lr{RS-DRL} اختصاص دارد. جزئیات الگوریتمیِ مکانیزم شکل‌دهی تصادفی پاداش، نحوه مدیریت حافظه بازپخش، و فرآیند آموزش عامل یادگیرنده به تفصیل بیان می‌گردد. همچنین نحوه اعتبارسنجی گزینه‌های وفقی پیش از اجرا توسط ابزار \lr{UPPAAL-SMC} برای تضمین پایداری سیستم مورد بررسی قرار می‌گیرد.

\textbf{فصل ۵: ارزیابی تجربی}

نحوه تخصیص محیط شبیه‌سازی، سناریوهای متعدد پیکربندی شبکه، و پروتکل‌های مقایسه‌ی جامع مدل \lr{RS-DRL} با استراتژی‌های رقیب در این فصل ارائه می‌گردد. ارائه‌ی شفافِ معیارهای کلیدی دستاوردها شامل عملکرد مجانبی، سرعت همگرایی و تحلیل‌های آماری در این فصل مستند شده‌اند.

\textbf{فصل ۶: بحث و نتیجه‌گیری}

در این فصل، به بررسی مصالحه‌ها، محدودیت‌ها و چالش‌های باقی‌مانده پرداخته شده و نتیجه‌گیری کلی پیرامون مزیت‌های کاربردی راهکار نوین ارائه می‌شود. همچنین مسیرهایی جهت بسط رویکرد فعلی در آینده پیشنهاد می‌گردد.
\newpage
